% generated by docs/audit/make_tables.py -- do not edit by hand
\noindent{\footnotesize\ttfamily lean/Pair9Shift.lean:5296}\par\smallskip
{\footnotesize\ttfamily\setlength{\parindent}{0pt}\raggedright
\mbox{}theorem~nonVacuity~:\\[-1pt]
\mbox{}~~~~Nonempty~AKR9.Pair9\\[-1pt]
\mbox{}~~~~$\land$~(points9~AKR9.certPair9).Nonempty\\[-1pt]
\mbox{}~~~~$\land$~(points9~AKR9.certPair9).Finite\\[-1pt]
\mbox{}~~~~$\land$~($\forall$~v~j~:~Fin~7,~AKR9.certPair9.V~(pP~v)~$\ne$~AKR9.certPair9.V~(pQ~j))\\[-1pt]
\mbox{}~~~~$\land$~AKR9.NoVertexOnOther~AKR9.certPair9\\[-1pt]
\mbox{}~~~~$\land$~(points9~AKR9.certPair9).ncard~=~38\\[-1pt]
\mbox{}~~~~$\land$~$\lnot$~(points9~AKR9.certPair9).ncard~$\le$~37\\[-1pt]
\mbox{}~~~~$\land$~(38~:~$\mathbb{N}$)~$\in$~attainedCoincFree\\[-1pt]
\mbox{}~~~~$\land$~(39~:~$\mathbb{N}$)~$\notin$~attainedCoincFree\\[-1pt]
\mbox{}~~~~$\land$~IsGreatest~attainedCoincFree~38~:=\\[-1pt]
\mbox{}~~$\langle$pair9\_nonempty,~certPair9\_points9\_nonempty,~certPair9\_finite,\\[-1pt]
\mbox{}~~~certPair9\_coincFree,~certPair9\_noVertexOnOther,~AKR9.certPair9\_ncard,\\[-1pt]
\mbox{}~~~certPair9\_not\_le\_37,~mem\_attainedCoincFree\_38,~not\_mem\_attainedCoincFree\_39,\\[-1pt]
\mbox{}~~~isGreatest\_38\_of\_coincFree$\rangle$\\[-1pt]
}
\par\medskip\noindent Ten clauses, in order:
\begin{enumerate}[leftmargin=2.1em,label=\textbf{(\arabic*)}]
\item \code{Pair9} is inhabited: the two \code{SimpleClosed} hypotheses of \S2 are satisfiable at all. Without this the whole theory could be about an empty type.
\item The witness's intersection set is non-empty.
\item It is \emph{finite} --- so \code{ncard} is the real cardinality. On an infinite set \code{ncard} returns $0$, which would make the bound $\le 38$ true and meaningless; this clause closes that door. (Clause~(6) already implies it, for exactly that reason; it is stated separately for the reader.)
\item The witness satisfies \code{hcoinc}: it shares no vertex. The hypothesis of the theorem is not contradictory.
\item The witness also satisfies the \emph{stronger} condition \code{NoVertexOnOther} --- no vertex anywhere on the other boundary.
\item Its count is exactly $38$: \textbf{the bound is attained}.
\item $\le 37$ is \emph{false} there: \textbf{the bound is not slack}. One less is provably wrong, so $38$ is the right constant and not an artefact of a lossy estimate. (Immediate from~(6); stated because it is the thing a reader wants to see said.)
\item $38$ is an attained value.
\item $39$ is \emph{not} attained.
\item $38$ is the greatest attained value.
\end{enumerate}
\par\medskip\noindent These are ten clauses, not ten logically independent facts: (6) is the sharpest of them and already implies (3) and (7). The content is not the number of clauses but the fact that \textbf{(2)--(7) all hold of the \emph{same} configuration}, \code{AKR9.certPair9}. Checking hypotheses one at a time would not rule out a theorem whose hypotheses are each satisfiable but jointly contradictory; a single shared witness does. That is precisely the failure mode an absent human reviewer would be looking for, and the declaration's own doc-comment says so.
\par\medskip\noindent A companion records that the maximum is \emph{unique} --- any $n$ that is a greatest attained value equals $38$:\par\smallskip
\noindent{\footnotesize\ttfamily lean/Pair9Shift.lean:5272}\par\smallskip
{\footnotesize\ttfamily\setlength{\parindent}{0pt}\raggedright
\mbox{}theorem~isGreatest\_attainedCoincFree\_unique~\{n~:~$\mathbb{N}$\}\\[-1pt]
\mbox{}~~~~(h~:~IsGreatest~attainedCoincFree~n)~:~n~=~38~:=\\[-1pt]
\mbox{}~~h.unique~isGreatest\_38\_of\_coincFree\\[-1pt]
}
